\subsection{Сила Лоренца*}

%%%%%%%%%%%%%%%%%%%%%%%%%%%%%%%%%%%%%%%%%%%%%%%%%%%%
%%%%%%%%%%%%%%%%%%%%%%%%%%%%%%%%%%%%%%%%%%%%%%%%%%%%
 \z{Электрон, ускоренный напряжением $200\ \text{В}$, движется в магнитном поле Земли, индукция которого $70\ \text{мкТл}$. Найдите радиус окружности, по которой движется электрон, если скорость его перпендикулярна магнитному полю Земли.}
%
{$R=\sqrt{\frac{2mV}{e}}\frac{1}{B} \approx 0.68\,\text{м}$}

%%%%%%%%%%%%%%%%%%%%%%%%%%%%%%%%%%%%%%%%%%%%%%%%%%%%
 \z{Как относятся радиусы траекторий двух электронов с кинетической энергией $K_{1}$ и $K_{2}$, если однородное магнитное поле перпендикулярно их скорости?}
%
{$\frac{r_1}{r_2}=\sqrt{\frac{K_1}{K_2}}$}

%%%%%%%%%%%%%%%%%%%%%%%%%%%%%%%%%%%%%%%%%%%%%%%%%%%%
 \z{Пространство разделено на две области плоскостью. В одной области создано магнитное поле индукции $B_{1}$, в другой – индукции $B_{2}$. Причём поля однородны и параллельны друг другу. С плоскости раздела перпендикулярно ей стартует электрон со скоростью $v$ в сторону области с индукцией поля $B_{1}$. Опишите дальнейшее движение электрона. Определите среднюю (дрейфовую) скорость перемещения электрона вдоль границы раздела магнитных полей, проницаемой для него.}
%
{$v_{\text{дрейф}}=\frac{2v(B_2-B_1)}{\pi (B_2+B_1)}$}

%%%%%%%%%%%%%%%%%%%%%%%%%%%%%%%%%%%%%%%%%%%%%%%%%%%%
 \z{В однородном магнитном поле с индукцией $B$ движется со скоростью $v$ частица массы $m$ с зарядом $q$, причём угол между векторами $B$ и $v$ равен $30\degree$. Найти радиус и шаг спирали, по которой движется частица.}
%
{$r=\frac{m v_{\perp}}{|q| B}=\frac{m v}{2|q| B};$~$d=v_{\parallel}T=\frac{2\pi m v\cos30\degree}{|q| B}=2\pi r\ctg 30\degree=2\pi r\sqrt{3}.$}

%%%%%%%%%%%%%%%%%%%%%%%%%%%%%%%%%%%%%%%%%%%%%%%%%%%%
 \z{Проводящее кольцо поместили в магнитное поле, перпендикулярное его плоскости. По кольцу циркулирует ток $I$. Если проволока кольца выдерживает на разрыв нагрузку $F$, то при какой индукции магнитного поля кольцо разорвётся? Радиус кольца $R$. Действием на кольцо магнитного поля, создаваемого током $I$, пренебречь.}
%
{$B_{\text{разрыв}}=\frac{F}{I R}$}

